% Abstract (ICML): ONE paragraph, "roughly 4-6 sentences" (ICML 2026 template /
% camera-ready instructions); 120-200 prose words [current]. A results
% sentence may stay commented out [some]. No numbers, citations or dataset
% names [all]. Math symbols are allowed [some] (e.g. $s^{\text{div}}$).
% Slots [venue: ICML]: (1) concessive context + challenge; (2) prior work named;
% (3) gap pair "(i) ...; and (ii) ..."; (4) proposal "we propose X, which ...";
% (5) "Specifically, ..."; (6) the training-signal sentence.
\begin{abstract}
Although <task> matters for many real-world applications, <model family> still struggle with it because they suffer from <degradation>.
A recent method, <Prior>, tackles this problem by <mechanism of prior work>.
Yet this <prior design> has two key weaknesses:
(i) <drawback 1, with its cause>;
and (ii) <drawback 2, with its consequence>.
To resolve both, we propose \textbf{X}, which adds <two components> to <workflow> for more <epithet> <task>.
Specifically, in X, <component A behaviour>, and <component B behaviour>.
We learn both behaviours with <two training signals> inside <training paradigm>, which reward <behaviour A> and <behaviour B>, respectively.
% Across <task family>, X improves on <Prior> by up to <gain>.   % results sentence kept commented [some]
\end{abstract}
